\chapter{KESIMPULAN DAN SARAN}

\section{Kesimpulan}

Berdasarkan hasil eksperimen dan pembahasan pada Bab~IV, penelitian ini
mengevaluasi penggunaan representasi laten sebagai media penalaran dan
komunikasi pada sistem multi-agen melalui perbandingan lima formulasi langkah
laten, yaitu \textit{raw}, \textit{soft}, \textit{sample}, \textit{Gumbel-Softmax},
dan \textit{Mixture of Inputs} (MoI). Evaluasi dilakukan pada dua lengan
eksperimen, yaitu lengan \textit{benchmark} untuk mengukur kemampuan penalaran
dan efisiensi komunikasi serta lengan faktor simbolik untuk menguji kemampuan
mempertahankan informasi berbentuk ekspresi terstruktur. Berdasarkan hasil
tersebut, kesimpulan penelitian dapat dirumuskan sebagai berikut.

\subsection{Formulasi Representasi Laten dalam Satu Kerangka Matematis}

Berbagai formulasi langkah laten yang digunakan dalam penelitian dapat
dinyatakan dalam kerangka umum
\[
z = \sum_{i=1}^{V} w_i W_{\mathrm{in}}[i],
\qquad
w\in\Delta^{V-1},
\]
untuk keluarga metode yang menghasilkan kombinasi konveks embedding token.
Dalam kerangka tersebut, perbedaan utama antara \textit{soft}, \textit{sample},
\textit{Gumbel-Softmax}, dan \textit{Mixture of Inputs} terletak pada aturan
pembentukan bobot ($w$), sedangkan persamaan \textit{raw} pada LatentMAS tidak
memiliki pembatasan bahwa koefisien harus tidak negatif dan berjumlah satu.
Dengan demikian, \textit{raw} berada di luar keluarga relaksasi diskret yang
dibandingkan.

Pada MoI, bobot dapat dinyatakan sebagai kombinasi distribusi prediksi dan
token hasil sampling. Dengan parameter ($\beta=1$), bobot MoI berada pada
keluarga kombinasi konveks antara representasi \textit{soft} dan representasi
hasil sampling. Gumbel-Softmax juga menghasilkan interpolasi antara distribusi
kontinu dan kecenderungan menuju token diskret, tetapi melalui manipulasi
logit dan parameter suhu.

Perbedaan geometris tersebut konsisten dengan hasil pengukuran representasi.
Kosinus antara langkah laten dan embedding token terdekat mencapai $0{,}927$ pada
\textit{soft}, $0{,}985$ pada Gumbel, dan $1{,}000$ pada \textit{sample}, sedangkan
\textit{raw} hanya mencapai $0{,}312$. Hasil ini menunjukkan bahwa anggota keluarga
relaksasi menghasilkan representasi yang secara geometris lebih dekat dengan
ruang embedding token dibandingkan persamaan \textit{raw}.

\subsection{Pengaruh Formulasi terhadap Penalaran dan Presisi Simbolik}

Hasil lengan \textit{benchmark} menunjukkan bahwa pengaruh formulasi langkah
laten tidak seragam untuk seluruh jenis tugas. Pada GSM8K dan ARC-Challenge,
kelima formulasi menghasilkan kinerja yang relatif berdekatan dan tidak
menunjukkan perbedaan yang cukup kuat untuk menyimpulkan bahwa salah satu
formulasi secara umum lebih unggul. Sebaliknya, pada HumanEval-Plus,
\textit{raw} menunjukkan penurunan kinerja yang jelas dengan akurasi ($0{,}42$),
sementara formulasi lainnya berada pada rentang ($0{,}69$)--($0{,}74$).

Pengujian berpasangan menunjukkan bahwa perbedaan signifikan yang bertahan
setelah koreksi multipisitas hanya melibatkan \textit{raw} pada HumanEval-Plus.
Tidak ditemukan perbedaan signifikan antaranggota keluarga relaksasi. Dengan
demikian, hasil penelitian lebih mendukung pengaruh \textit{keanggotaan
keluarga formulasi} daripada keunggulan salah satu metode relaksasi tertentu.

Pengaruh tersebut juga bergantung pada tuntutan presisi simbolik tugas.
Perbedaan kerusakan antara HumanEval-Plus dan GSM8K maupun antara HumanEval-Plus
dan ARC-Challenge signifikan, sedangkan GSM8K dan ARC-Challenge tidak berbeda
secara meyakinkan. Oleh karena itu, data tidak mendukung urutan tiga tingkat
tuntutan simbolik yang dirumuskan sebelumnya, tetapi mendukung pembedaan dua
kelompok, yaitu tugas pembangkitan kode dan kedua tugas penalaran lainnya.

Temuan tersebut diperkuat oleh pengamatan terhadap fidelitas simbolik.
Kejadian korupsi token lebih banyak ditemukan pada \textit{raw}, termasuk
kemunculan karakter yang tidak diharapkan serta kerusakan pada tingkat
sub-kata. Pada tugas yang membutuhkan keluaran simbolik secara presisi,
kerusakan semacam ini lebih mudah menyebabkan kegagalan keseluruhan dibanding
pada tugas dengan jawaban pendek.

Dengan demikian, kelemahan yang ditemukan bukan menunjukkan bahwa komunikasi
melalui KV-\textit{cache} secara inheren kehilangan informasi. Bukti yang diperoleh
justru menunjukkan bahwa masalah utama berada pada formulasi yang digunakan
untuk membentuk representasi laten sebelum representasi tersebut dimasukkan ke
dalam KV-\textit{cache}.

\subsection{Efisiensi Komunikasi Laten}

Komunikasi antar-agen melalui KV-\textit{cache} menunjukkan pengurangan biaya yang
substansial dibandingkan komunikasi berbasis teks. Penghematan token keluaran
mencapai $75{,}0\%$--$87{,}8\%$, sedangkan percepatan waktu inferensi mencapai
$1{,}8\times$--$6{,}1\times$. Rentang penghematan token tersebut mencakup
rentang $70{,}8\%$--$83{,}7\%$ yang dilaporkan oleh \textit{LatentMAS}, sehingga temuan
efisiensi kerangka komunikasi laten dapat direplikasi pada implementasi,
perangkat keras, dan konfigurasi eksperimen penelitian ini.

Penghematan terutama berasal dari berkurangnya keluaran teks antar-agen.
Pada medium teks, setiap agen perantara harus menghasilkan keluaran tekstual
yang kemudian diproses oleh agen berikutnya. Pada medium KV-\textit{cache}, representasi
internal diteruskan secara langsung sehingga proses verbalisasi penalaran
antar-agen dapat dihilangkan.

Namun, efisiensi tersebut tidak secara otomatis berarti peningkatan akurasi.
Pada ketiga tugas, medium KV dengan formulasi keluarga relaksasi menghasilkan
kinerja yang secara umum sebanding dengan medium teks, tanpa menunjukkan
keunggulan akurasi yang konsisten. Oleh karena itu, temuan penelitian lebih
tepat dirumuskan sebagai
\[
\boxed{
\text{akurasi sebanding}
\quad+\quad
\text{biaya komunikasi lebih rendah}
}
\]
daripada sebagai keunggulan akurasi komunikasi laten terhadap komunikasi teks.

Efisiensi tersebut juga bergantung pada formulasi representasi laten. Pada
HumanEval-Plus, penggunaan \textit{raw} menghasilkan akurasi ($0{,}42$) pada
medium KV, jauh di bawah medium teks sebesar ($0{,}76$) dan bahkan di bawah
konfigurasi agen tunggal sebesar ($0{,}72$). Dengan demikian, efisiensi tidak
dapat dipandang sebagai sifat medium KV-\textit{cache} saja, tetapi sebagai sifat
gabungan antara medium komunikasi dan formulasi representasi laten yang
mengisinya.

\subsection{Kemampuan Mempertahankan Muatan Simbolik pada Generasi Faktor}

Lengan faktor simbolik memberikan pengujian tambahan terhadap kemampuan
representasi laten dalam membawa informasi yang harus dipertahankan dalam
bentuk simbolik dan terstruktur. Pada lengan ini, \textit{raw} menunjukkan
tingkat kegagalan keluaran yang paling tinggi: 14 dari 18 panggilan tidak
menghasilkan ekspresi yang dapat diurai, atau sebesar $78\%$. Sebaliknya,
anggota keluarga relaksasi menghasilkan keluaran yang lebih dapat diandalkan,
dengan tingkat kegagalan berada pada kisaran $29\%$--$40\%$ untuk metode
yang mengalami kegagalan.

Perbedaan tersebut berlanjut pada jumlah ekspresi yang dapat dievaluasi.
\textit{Raw} menghasilkan 11 ekspresi dengan hanya 3 ekspresi yang dapat
dievaluasi, sedangkan anggota keluarga relaksasi menghasilkan 22--33 ekspresi
dengan 21--26 ekspresi yang dapat dievaluasi. Rerata
($\lvert\mathrm{RankIC}\rvert$) pada ekspresi yang dapat dievaluasi berada
pada kisaran ($0{,}0162$)--($0{,}0185$) untuk medium teks dan anggota keluarga
relaksasi, sedangkan \textit{raw} menghasilkan ($0{,}0050$) berdasarkan hanya
tiga ekspresi.

Hasil tersebut memberikan bukti tambahan bahwa kemampuan kanal laten dalam
mempertahankan muatan simbolik sangat dipengaruhi oleh formulasi representasi
yang digunakan. Akan tetapi, karena jumlah jalan pada lengan faktor relatif
kecil dan satuan pengamatannya berbeda dari lengan \textit{benchmark}, hasil
ini diperlakukan sebagai bukti pendukung dan tidak digunakan sebagai dasar
pengujian statistik formal.

Nilai \textit{backtest} pada lengan faktor tidak ditafsirkan sebagai prediksi
keuntungan investasi. Penggunaannya terbatas pada pengujian apakah ekspresi
simbolik yang berhasil dibawa melalui sistem dapat dieksekusi dan menghasilkan
sinyal yang dapat dievaluasi secara kuantitatif.

\subsection{Kesimpulan Umum}

Secara keseluruhan, penelitian menunjukkan bahwa representasi laten dapat
digunakan sebagai media komunikasi antar-agen yang jauh lebih efisien daripada
komunikasi berbasis teks tanpa harus mengorbankan kinerja penalaran pada tugas
yang diuji, tetapi keberhasilan tersebut tidak hanya ditentukan oleh medium
komunikasinya. Formulasi pembentukan representasi laten menjadi komponen
penting yang menentukan apakah informasi yang diteruskan tetap berada pada
bentuk yang dapat diproses secara efektif oleh model.

Keluarga formulasi yang menghasilkan kombinasi konveks embedding token,
yaitu \textit{soft}, \textit{sample}, \textit{Gumbel-Softmax}, dan MoI,
menunjukkan perilaku yang relatif konsisten pada ketiga \textit{benchmark}
serta lebih mampu mempertahankan informasi simbolik dibandingkan formulasi
\textit{raw} pada tugas dengan tuntutan presisi simbolik tinggi. Sebaliknya,
persamaan \textit{raw} yang digunakan dalam LatentMAS menunjukkan kelemahan
terutama ketika keluaran menuntut ketepatan identitas simbol.

Dengan demikian, temuan utama penelitian bukan bahwa satu metode relaksasi
tertentu merupakan formulasi terbaik, melainkan bahwa \textbf{pemilihan
ruang representasi tempat langkah laten dibentuk merupakan faktor penting
dalam keberhasilan latent collaboration}. Medium KV-\textit{cache} menyediakan jalur
komunikasi yang efisien, sedangkan formulasi langkah laten menentukan apakah
informasi yang dimasukkan ke jalur tersebut tetap dapat digunakan secara
efektif oleh agen berikutnya.

\section{Saran}

Beberapa arah lanjutan dapat ditempuh berdasarkan batas yang dinyatakan pada
Bab~IV.

\begin{enumerate}
    \item \textit{Pengulangan lintas seed dan perbesaran sampel.}
          Setiap sel pada penelitian ini dijalankan satu kali sehingga ragam
          antar-\textit{seed} tidak terukur. Pengulangan dengan beberapa
          \textit{seed} akan memisahkan efek metode dari ragam pembangkitan,
          terutama untuk selisih kecil pada GSM8K dan ARC-Challenge yang saat
          ini dilaporkan sebagai tidak signifikan.

    \item \textit{Pengujian lintas backbone.} Penelitian ini
          mengidentifikasi keterikatan matriks embedding (\textit{tied
          embeddings}) sebagai variabel yang menentukan apakah matriks
          \textit{realignment} bekerja atau tidak. Variabel tersebut dapat
          diuji langsung dengan membandingkan model yang matriksnya terikat dan
          tidak terikat pada ukuran yang sebanding.

    \item \textit{Pengujian medium gabungan pada ukuran sampel memadai.}
          Medium gabungan KV dan teks belum diuji sebagai kondisi penuh. Selain
          itu, penyebab hilangnya prefiks keluaran pada lengan faktor perlu
          ditelusuri sampai ke akar implementasinya, karena pada penelitian ini
          bentuk kegagalannya dipulihkan tanpa dijelaskan.

    \item \textit{Pengujian langsung hipotesis mekanistik.} Kerangka pada
          Bab~III memungkinkan pembentukan persamaan antara yang berada pada
          jarak terkendali dari lambung konveks, misalnya melalui interpolasi
          antara persamaan resmi dan bobot \textit{soft}. Apabila akurasi
          bergerak secara monoton mengikuti jarak tersebut, hubungan antara
          geometri representasi dan kinerja tugas menjadi terbukti secara
          kausal, bukan sekadar konsisten.

    \item \textit{Perluasan topologi dan domain.} Rantai agen pada penelitian
          ini bersifat sekuensial dengan panjang tetap. Topologi hierarkis,
          jumlah agen yang berbeda, dan lengan faktor simbolik yang dijalankan
          sebagai kondisi penuh --- termasuk lapisan evolusi yang sengaja tidak
          diaktifkan di sini --- merupakan perluasan yang wajar.

    \item \textit{Penilaian ekonomi yang lebih ketat pada domain faktor.}
          Metrik \textit{backtest} pada penelitian ini dihitung tanpa
          \textit{slippage}, tanpa batas likuiditas, dan dengan penyeimbangan
          ulang harian penuh. Penilaian yang memperhitungkan biaya transaksi
          dan kendala perdagangan diperlukan sebelum angka tersebut dapat
          dibaca sebagai kinerja investasi.
\end{enumerate}
